\documentclass[10pt,a4paper]{amsart}
\usepackage[margin=19mm]{geometry}
\usepackage{amsmath,amssymb,mathtools}
\usepackage[scr=rsfs,cal=euler]{mathalfa}
\usepackage{xfrac}
\usepackage[unicode,hidelinks,bookmarks=false]{hyperref}
\newcommand{\ie}{i.e.\ }
\newcommand{\cf}{cf.\ }
\newcommand{\resp}{respectively\ }
\newcommand{\Subsection}{\subsection}
\providecommand{\nobreakdash}{\nobreak\hbox{-}}
\setlength{\emergencystretch}{2em}
\multlinegap0pt
\usepackage{bm}
\usepackage{bbm}
\usepackage{dsfont}
\usepackage{xspace}
\usepackage{cancel}
\usepackage{mathtools}
\usepackage{cleveref}
\usepackage{amsthm}
\makeatletter
\@ifundefined{lemma}{\newtheorem{lemma}{Lemma}}{}
\@ifundefined{proposition}{\newtheorem{proposition}{Proposition}}{}
\makeatother
\newcommand{\E}{\mathbb E}
\newcommand{\Prob}{\mathbb P}
\newcommand{\Z}{\mathbb Z}
\newcommand{\floor}[1]{\left\lfloor #1\right\rfloor}
\title[Long Lattice Paths with No Three Collinear Vertices]{Long Lattice Paths with No Three Collinear Vertices}
\author{Samuel Korsky}
\date{Source: arXiv:2608.07906v1 (CC BY 4.0)}
\begin{document}
\maketitle

\noindent\emph{Source and licence.} Long Lattice Paths with No Three Collinear Vertices. Version arXiv:2608.07906v1 (CC BY 4.0), distributed under the Creative Commons Attribution 4.0 International licence, as recorded in the arXiv licence metadata for this exact version and shown on the abstract page https://arxiv.org/abs/2608.07906v1. Full rights notice: licenses/arxiv-2608.07906-CC-BY-4.0.txt.\par\smallskip

\noindent\textit{Editorial scope.} Counted unit: the Proposition whose source label is \texttt{prop:amplification} in the article (source0.tex lines 408--417, source label \texttt{prop:amplification}), delivered together with the article's own complete original proof of it (source0.tex lines 419--572, one \texttt{proof} environment of 154 lines). No mathematics is written, repaired or re-derived: statement and proof are the article's own text line for line. Required context retained uncounted: lem:periodic-reduction (lines 155--161); lem:weighted-uniform-sums (lines 343--359). Every definition, display and label of those lines is retained unchanged. Omitted from this package and not redistributed: 1--154, 162--342, 360--407, 418--418, 573--824. Nothing inside a retained excerpt is omitted or altered: every retained body is delivered exactly as the article's lines are written, so no omission receipt of any kind accompanies this package. Citations: the retained excerpts carry no citation command at all, so no bibliography entry of the article is transcribed and no reference is redistributed. Reference adaptation: none was needed and none was made; every reference inside the delivered unit resolves inside it, so no unresolved reference or citation is produced. Printed slips kept literal: no printed word, symbol, hypothesis or number of the counted statement or of its proof was altered. Layout and class: the article's own class is \texttt{article} with packages including geometry, amsmath, amssymb, amsthm, mathtools, enumitem, xcolor, hyperref, cleveref; the wrapper is the lane's pinned \texttt{amsart} editorial wrapper at 10pt on A4, which restates the theorem-like environments the retained text needs and adds the article's own macro definitions that the retained text needs (\texttt{\textbackslash E}, \texttt{\textbackslash Prob}, \texttt{\textbackslash Z}, \texttt{\textbackslash floor}), with extra wrapper packages (bm, bbm, dsfont, xspace, cancel, mathtools, cleveref). The delivered numbering is the unit's own. Assets: no retained span contains a figure environment, a graphics-inclusion command, a PDF-page inclusion, a table or a TikZ picture; no image file is copied into this package, the package declares no asset dependency and no image file is redistributed with it. Licence: version arXiv:2608.07906v1 of this article is distributed under the Creative Commons Attribution 4.0 International licence, as recorded in the arXiv licence metadata for this exact version and shown on the abstract page https://arxiv.org/abs/2608.07906v1. A full rights notice is delivered at licenses/arxiv-2608.07906-CC-BY-4.0.txt. Characters: every retained excerpt is pure ASCII, so the delivered engine, XeLaTeX with its default Latin Modern text font, reports no missing character.
\begin{lemma}[Periodic reduction]\label{lem:periodic-reduction}
Let \(X\) be a 3-free word of length \(n\), let \(\#\) be a new letter, and let
\[
        W=(X\#)^\infty
\]
have period \(M=n+1\). Two adjacent nonempty intervals of \(W\) form a bad pair if and only if both lengths are multiples of \(M\).
\end{lemma}

\begin{lemma}[Weighted point-mass bound]\label{lem:weighted-uniform-sums}
Let $h,q,r,a,b\ge1$. Let $U_1,\ldots,U_q,V_1,\ldots,V_r$ be independent and uniform on $\{-h,\ldots,h\}$, and put
\[
        S_q=\sum_{j=1}^qU_j,
        \qquad
        T_r=\sum_{j=1}^rV_j.
\]
If $D=\gcd(a,b)$, then
\begin{equation}\label{eq:weighted-uniform-sums}
 \sup_{z\in\Z}\Prob(bS_q-aT_r=z)
 \le C\left(
        \frac{1}{h^2\sqrt{qr}}
        +\frac{D}{h\sqrt{a^2r+b^2q}}
 \right)
\end{equation}
for an absolute constant $C$.
\end{lemma}

\begin{proposition}[Amplification]\label{prop:amplification}
Fix an integer $p\ge3$. There are constants $c_p,C_p>0$ with the following property. Let $(X,\sigma,\rho)$ be a marked 3-free word of length $n$ using \(r\) letters. If
\begin{equation}\label{eq:amplifier-size-hypothesis}
        n\ge C_p\rho^{-2},
\end{equation}
then there is a marked 3-free word \((Y,\tau,c_p\rho)\) using at most \(r+p+1\) letters such that
\begin{equation}\label{eq:amplifier-conclusion}
        c_p\rho^p n^p\le |Y|\le n^p.
\end{equation}
\end{proposition}

\begin{proof}
During the proof, \(c_p\) may be decreased and \(C_p\) increased; only finitely many such changes are made, so their final values are fixed. We may also assume \(0<c_p\le1\). Put
\[
        s=|X|_\sigma\ge\rho n,
        \qquad
        M=n+1,
\]
and form \(W=(X\#)^\infty\), where \(\#\) is a new letter. The size hypothesis implies
\[
        s\ge C_p\rho^{-1}
        \quad\text{and}\quad
        \frac{s^p}{M}
        \ge \frac{\rho^pn^{p-1}}{2}
        \ge \frac{C_p^{p-1}}{2}\cdot\rho^{2-p}
        \ge \frac{C_p^{p-1}}{2}.
\]
Thus, by increasing \(C_p\), we may assume that every later lower bound depending only on \(p\) is satisfied.

\emph{Random split.}
Choose a fixed integer \(K=K(p)\ge4\). Number the occurrences of \(\sigma\) in \(W\) as \(1,2,\ldots\), and let each \(\sigma\)-block consist of
\[
        g=\floor{s/K}
\]
consecutive occurrences in this ordering. Put
\[
        u=\floor{g/(p+1)},
        \qquad
        h=\floor{u/(10p)}.
\]
For each \(\sigma\)-block \(B\), independently choose
\[
        U_{B,1},\ldots,U_{B,p}
\]
uniformly from \(\{-h,\ldots,h\}\). Replace the \(g\) occurrences of \(\sigma\) in \(B\) by new letters \(\sigma_1,\ldots,\sigma_{p+1}\) with respective counts
\begin{equation}\label{eq:block-composition}
        u+U_{B,1},\ldots,u+U_{B,p},
        \qquad
        g-pu-\sum_{i=1}^pU_{B,i}.
\end{equation}
Assign the first prescribed number of \(\sigma\)-positions in \(B\) to \(\sigma_1\), the next prescribed number to \(\sigma_2\), and so on.

Write \(g=(p+1)u+r_0\), where \(0\le r_0\le p\). The first \(p\) counts in \eqref{eq:block-composition} are at least \(u-h\), and the last is at least \(u-ph\). For sufficiently large \(u\), all are therefore at least \(\delta_pg\) for some \(\delta_p>0\). Moreover,
\[
        g\ge \frac{s}{2K},\qquad
        u\ge \frac{g}{2(p+1)},\qquad
        h\ge \frac{u}{20p},
\]
so
\begin{equation}\label{eq:h-comparable-s}
        h\ge c_ps.
\end{equation}

Choose \(0<\varepsilon_p\le1\) sufficiently small and set
\begin{equation}\label{eq:amplifier-length}
        N=M\floor{\frac{\varepsilon_ps^p}{M}}.
\end{equation}
By the preceding lower bound on \(s^p/M\), we may assume that the quantity inside the floor is at least \(2\). Since \(\lfloor x\rfloor\ge x/2\) for \(x\ge2\),
\begin{equation}\label{eq:N-bounds}
        \frac{\varepsilon_p}{2}\cdot s^p\le N\le \varepsilon_ps^p\le n^p.
\end{equation}
Let \(Y\) be the first \(N\) letters of the split word.

\emph{Possible bad pairs.}
Any bad pair in \(Y\) projects, after merging \(\sigma_1,\ldots,\sigma_{p+1}\) back to \(\sigma\), to a bad pair in \(W\). By \cref{lem:periodic-reduction}, its two lengths are \(aM\) and \(bM\) for positive integers \(a,b\). For each fixed pair \((a,b)\), there are at most \(N\) possible starting positions.

Fix one possible bad pair, with adjacent intervals \(I\) and \(J\) of lengths \(aM\) and \(bM\). Each interval of length \(aM\) in \(W\) contains exactly \(as\) occurrences of \(\sigma\), regardless of its starting position. Let \(q_I\) and \(q_J\) be the numbers of complete \(\sigma\)-blocks contained in \(I\) and \(J\), respectively. At most two \(\sigma\)-blocks meet an interval without being contained in it. Since \(g\le s/K\),
\begin{equation}\label{eq:complete-block-counts}
        q_I\ge \frac{as}{g}-2\ge (K-2)a,
        \qquad
        q_J\ge \frac{bs}{g}-2\ge (K-2)b.
\end{equation}

Condition on every random variable except those attached to complete \(\sigma\)-blocks contained in \(I\) or \(J\). If the pair remains bad after the split, equality of the normalized counts of \(\sigma_i\), for each \(1\le i\le p\), imposes an equation
\begin{equation}\label{eq:coordinate-collision}
        b\sum_{B\subset I}U_{B,i}
        -a\sum_{B\subset J}U_{B,i}
        =z_i,
\end{equation}
where \(z_i\) is determined by the conditioned variables, the baseline counts, and the partially intersected blocks. The \(p\) equations use independent families of random variables. The last replacement letter gives no further condition, because its count is determined by the first \(p\) counts and the total number of replaced occurrences.

Let \(D=\gcd(a,b)\). From \eqref{eq:h-comparable-s} and \eqref{eq:complete-block-counts},
\[
 h^2\sqrt{q_Iq_J}\ge c_ps^2\sqrt{ab},
 \qquad
 a^2q_J+b^2q_I\ge c_pab(a+b).
\]
Applying \cref{lem:weighted-uniform-sums} therefore shows that the conditional probability of one equation in \eqref{eq:coordinate-collision} is at most
\begin{equation}\label{eq:one-coordinate-survival}
        C_p\left(
        \frac{1}{s^2\sqrt{ab}}
        +\frac{D}{s\sqrt{ab(a+b)}}
        \right).
\end{equation}
Set
\[
 A=\frac{1}{s^2\sqrt{ab}},
 \qquad
 B=\frac{D}{s\sqrt{ab(a+b)}}.
\]
Conditional independence across the \(p\) coordinates gives a bound \(C_p(A+B)^p\). Since
\[
        (A+B)^p\le 2^{p-1}(A^p+B^p),
\]
and \(p\) is fixed, the factor \(2^{p-1}\) may be absorbed into \(C_p\). Hence
\begin{equation}\label{eq:candidate-survival}
 \Prob(\text{the pair remains bad})
 \le C_p\left(
        \frac{1}{s^{2p}(ab)^{p/2}}
        +\frac{D^p}{s^p[ab(a+b)]^{p/2}}
 \right).
\end{equation}

\emph{Expected number of bad pairs.}
The two sums
\begin{equation}\label{eq:summable-series}
 \sum_{a,b\ge1}\frac{1}{(ab)^{p/2}}
 \quad\text{and}\quad
 \sum_{a,b\ge1}
 \frac{\gcd(a,b)^p}{[ab(a+b)]^{p/2}}
\end{equation}
are finite for $p\ge3$. For the second, write $a=Du$ and $b=Dv$, where $D=\gcd(a,b)$, and then discard the coprimality condition:
\[
 \sum_{a,b\ge1}
 \frac{\gcd(a,b)^p}{[ab(a+b)]^{p/2}}
 \le
 \sum_{D\ge1}D^{-p/2}
 \sum_{u,v\ge1}[uv(u+v)]^{-p/2}<\infty.
\]
Indeed, $u+v\ge2\sqrt{uv}$, so the inner summand is at most a constant times $(uv)^{-3p/4}$.

Let \(Z\) be the number of bad pairs in \(Y\). Summing \eqref{eq:candidate-survival} over at most \(N\) starting positions for each \(a,b\), and using the convergence of the two displayed series, gives
\begin{equation}\label{eq:expected-bad-pairs}
        \E \left[Z\right]
        \le C_pN(s^{-p}+s^{-2p})
        \le C_p\varepsilon_p(1+s^{-p}).
\end{equation}
Choose \(\varepsilon_p\) so that the last quantity is less than \(1\). Since \(Z\) is a nonnegative integer, some outcome has \(Z=0\), and the corresponding word \(Y\) is 3-free.

\emph{Retaining a frequent letter.}
Since $M\mid N$, write $N=\ell M$. Before splitting, this prefix contains exactly \(\ell s=Ns/M\) occurrences of \(\sigma\). At most two \(\sigma\)-blocks meeting the prefix are incomplete, and every complete block contains at least $\delta_pg$ occurrences of $\sigma_1$. Therefore
\begin{equation}\label{eq:output-density}
        |Y|_{\sigma_1}
        \ge \delta_p(\ell s-2g).
\end{equation}
The floor in \eqref{eq:amplifier-length} is at least \(2\), so \(\ell\ge2\). Since \(K\ge4\) and \(g\le s/K\), this implies \(\ell s\ge4g\). Also \(M\le2n\) and \(s\ge\rho n\), so \eqref{eq:output-density} gives
\[
        |Y|_{\sigma_1}
        \ge c_p\cdot\frac{Ns}{M}
        \ge c_p\rho N.
\]
Thus we may take $\tau=\sigma_1$.

Finally, replacing \(\sigma\) by \(p+1\) letters increases the alphabet size by \(p\), and the separator adds one more letter. Thus the net increase is \(p+1\). The lower bound in \eqref{eq:amplifier-conclusion} follows from \eqref{eq:N-bounds} and \(s\ge\rho n\), after decreasing the fixed \(c_p\) if necessary.
\end{proof}

\end{document}
